\documentclass[11pt]{article}

% Common class-neutral preamble for standalone notes under manuscript/notes/.
% Note entry points all live exactly two directories below manuscript/.

\usepackage[margin=1in]{geometry}
\usepackage{amsmath,amssymb,amsthm,mathtools,bm}
\usepackage{algorithm}
\usepackage{algorithmic}
\usepackage{booktabs,tabularx,multirow,array,longtable}
\usepackage{enumitem}
\usepackage{microtype}
\usepackage[numbers,sort&compress]{natbib}
\usepackage{xcolor}
\usepackage[colorlinks=true,citecolor=blue,linkcolor=black,urlcolor=blue]{hyperref}
\usepackage[capitalize,nameinlink,noabbrev]{cleveref}

\newtheorem{theorem}{Theorem}[section]
\newtheorem{lemma}[theorem]{Lemma}
\newtheorem{proposition}[theorem]{Proposition}
\newtheorem{corollary}[theorem]{Corollary}
\newtheorem{conjecture}[theorem]{Conjecture}
\theoremstyle{definition}
\newtheorem{definition}[theorem]{Definition}
\newtheorem{assumption}[theorem]{Assumption}
\newtheorem{problem}[theorem]{Problem}
\newtheorem{observation}[theorem]{Observation}
\theoremstyle{remark}
\newtheorem{remark}[theorem]{Remark}
\theoremstyle{plain}

% Canonical mathematical notation for the active manuscript.
%
% This file preserves the recurring notation style used in the author's
% NeurIPS 2024 and 2025 papers while excluding unrelated tensor and
% machine-learning boilerplate from those archived command collections.
%
% Prefix convention:
%   \v...  bold vectors
%   \m...  bold matrices
%   \g...  calligraphic sets, graphs, and families
%   \s...  blackboard-bold spaces and number systems

\newcommand{\method}{\textsc{HybridLocalSolver}}

% Font and scalar shorthands retained from the archived manuscripts.
\newcommand{\mc}{\mathcal}
\newcommand{\R}{\mathbb{R}}
\newcommand{\E}{\mathbb{E}}
\newcommand{\G}{\mathbb{G}}
\newcommand{\N}{\mathcal{N}}
% Accuracy namespaces are semantic: do not add a bare \eps alias.
% Degree-normalized residual tolerance of classical APPR is kept distinct from
% the objective-gap and proximal fixed-point tolerances of the RPPR sections.
\newcommand{\epsappr}{\varepsilon_{\mathrm{appr}}}
\newcommand{\epsobj}{\varepsilon_{\mathrm{obj}}}
\newcommand{\epspg}{\varepsilon_{\mathrm{pg}}}
\newcommand{\epsppr}{\varepsilon_{\mathrm{ppr}}}
\newcommand{\epsin}{\varepsilon_{\mathrm{in}}}
\newcommand{\epskkt}{\varepsilon_{\mathrm{kkt}}}
\newcommand{\epsburn}{\varepsilon_{\mathrm{burn}}}
\newcommand{\epssol}{\varepsilon_{\mathrm{sol}}}
\newcommand{\epspprinst}[1]{\varepsilon_{\mathrm{ppr},#1}}
\newcommand{\trans}{\mathsf{T}}
\newcommand{\grad}{\nabla}

% Delimiters and generic helpers.
% Norms and inner products use fixed-size delimiters by default.
% Use an optional size (e.g. \norm[\big]{...}) or * for automatic sizing.
\DeclarePairedDelimiter{\norm}{\lVert}{\rVert}
\newcommand{\abs}[1]{\left\lvert #1 \right\rvert}
\DeclarePairedDelimiterX{\inner}[2]{\langle}{\rangle}{#1, #2}
\newcommand{\ip}[2]{\inner{#1}{#2}}
\newcommand{\card}[1]{\left\lvert #1 \right\rvert}
\newcommand{\set}[1]{\left\{#1\right\}}
\newcommand{\ceil}[1]{\left\lceil #1 \right\rceil}
\newcommand{\floor}[1]{\left\lfloor #1 \right\rfloor}
\newcommand{\comp}[1]{\overline{#1}}
\newcommand{\conj}[1]{\overline{#1}}
\newcommand{\mean}[1]{\overline{#1}}
\newcommand{\bigO}{\mathcal{O}}
\newcommand{\softO}{\widetilde{\mathcal{O}}}
\newcommand{\softOmega}{\widetilde{\Omega}}
\newcommand{\pos}[1]{\left[#1\right]_{+}}

% Bold lowercase vectors.
\newcommand{\va}{\bm{a}}
\newcommand{\vb}{\bm{b}}
\newcommand{\vc}{\bm{c}}
\newcommand{\vd}{\bm{d}}
\newcommand{\ve}{\bm{e}}
\newcommand{\vf}{\bm{f}}
\newcommand{\vg}{\bm{g}}
\newcommand{\vh}{\bm{h}}
\newcommand{\vi}{\bm{i}}
\newcommand{\vj}{\bm{j}}
\newcommand{\vk}{\bm{k}}
\newcommand{\vl}{\bm{l}}
\newcommand{\vm}{\bm{m}}
\newcommand{\vn}{\bm{n}}
\newcommand{\vo}{\bm{o}}
\newcommand{\vp}{\bm{p}}
\newcommand{\vq}{\bm{q}}
\newcommand{\vr}{\bm{r}}
\newcommand{\vs}{\bm{s}}
\newcommand{\vt}{\bm{t}}
\newcommand{\vu}{\bm{u}}
\newcommand{\vv}{\bm{v}}
\newcommand{\vw}{\bm{w}}
\newcommand{\vx}{\bm{x}}
\newcommand{\vy}{\bm{y}}
\newcommand{\vz}{\bm{z}}

% Bold Greek vectors and special vectors.
\newcommand{\valpha}{\bm{\alpha}}
\newcommand{\vbeta}{\bm{\beta}}
\newcommand{\vdelta}{\bm{\delta}}
\newcommand{\vepsilon}{\bm{\epsilon}}
\newcommand{\veta}{\bm{\eta}}
\newcommand{\vgamma}{\bm{\gamma}}
\newcommand{\vlambda}{\bm{\lambda}}
\newcommand{\vmu}{\bm{\mu}}
\newcommand{\vomega}{\bm{\omega}}
\newcommand{\vphi}{\bm{\phi}}
\newcommand{\vpi}{\bm{\pi}}
\newcommand{\vpsi}{\bm{\psi}}
\newcommand{\vrho}{\bm{\rho}}
\newcommand{\vsigma}{\bm{\sigma}}
\newcommand{\vtheta}{\bm{\theta}}
\newcommand{\vxi}{\bm{\xi}}
\newcommand{\vell}{\bm{\ell}}
\newcommand{\vDelta}{\bm{\Delta}}
\newcommand{\vGamma}{\bm{\Gamma}}
\newcommand{\vLambda}{\bm{\Lambda}}
\newcommand{\vPhi}{\bm{\Phi}}
\newcommand{\vSigma}{\bm{\Sigma}}
\newcommand{\vzero}{\bm{0}}
\newcommand{\vone}{\bm{1}}
\newcommand{\one}{\mathbf{1}}
\newcommand{\zero}{\vzero}
\newcommand{\eunit}[1]{\bm{e}_{#1}}
\newcommand{\vDeltax}{\bm{\Delta}\bm{x}}

% Bold uppercase matrices.
\newcommand{\mA}{\bm{A}}
\newcommand{\mB}{\bm{B}}
\newcommand{\mC}{\bm{C}}
\newcommand{\mD}{\bm{D}}
\newcommand{\mE}{\bm{E}}
\newcommand{\mF}{\bm{F}}
\newcommand{\mG}{\bm{G}}
\newcommand{\mH}{\bm{H}}
\newcommand{\mI}{\bm{I}}
\newcommand{\mJ}{\bm{J}}
\newcommand{\mK}{\bm{K}}
\newcommand{\mL}{\bm{L}}
\newcommand{\mM}{\bm{M}}
\newcommand{\mN}{\bm{N}}
\newcommand{\mO}{\bm{O}}
\newcommand{\mP}{\bm{P}}
\newcommand{\mQ}{\bm{Q}}
\newcommand{\mR}{\bm{R}}
\newcommand{\mS}{\bm{S}}
\newcommand{\mT}{\bm{T}}
\newcommand{\mU}{\bm{U}}
\newcommand{\mV}{\bm{V}}
\newcommand{\mW}{\bm{W}}
\newcommand{\mX}{\bm{X}}
\newcommand{\mY}{\bm{Y}}
\newcommand{\mZ}{\bm{Z}}

% Bold Greek matrices retained from the graph papers.
\newcommand{\mBeta}{\bm{\beta}}
\newcommand{\mDelta}{\bm{\Delta}}
\newcommand{\mGamma}{\bm{\Gamma}}
\newcommand{\mLambda}{\bm{\Lambda}}
\newcommand{\mOmega}{\bm{\Omega}}
\newcommand{\mPhi}{\bm{\Phi}}
\newcommand{\mPi}{\bm{\Pi}}
\newcommand{\mSigma}{\bm{\Sigma}}
\newcommand{\mTheta}{\bm{\Theta}}

% Calligraphic symbols.
\newcommand{\gA}{\mathcal{A}}
\newcommand{\gB}{\mathcal{B}}
\newcommand{\gC}{\mathcal{C}}
\newcommand{\gD}{\mathcal{D}}
\newcommand{\gE}{\mathcal{E}}
\newcommand{\gF}{\mathcal{F}}
\newcommand{\gG}{\mathcal{G}}
\newcommand{\gH}{\mathcal{H}}
\newcommand{\gI}{\mathcal{I}}
\newcommand{\gJ}{\mathcal{J}}
\newcommand{\gK}{\mathcal{K}}
\newcommand{\gL}{\mathcal{L}}
\newcommand{\gM}{\mathcal{M}}
\newcommand{\gN}{\mathcal{N}}
\newcommand{\gO}{\mathcal{O}}
\newcommand{\gP}{\mathcal{P}}
\newcommand{\gQ}{\mathcal{Q}}
\newcommand{\gR}{\mathcal{R}}
\newcommand{\gS}{\mathcal{S}}
\newcommand{\gT}{\mathcal{T}}
\newcommand{\gU}{\mathcal{U}}
\newcommand{\gV}{\mathcal{V}}
\newcommand{\gW}{\mathcal{W}}
\newcommand{\gX}{\mathcal{X}}
\newcommand{\gY}{\mathcal{Y}}
\newcommand{\gZ}{\mathcal{Z}}

% Blackboard-bold symbols.
\newcommand{\sA}{\mathbb{A}}
\newcommand{\sB}{\mathbb{B}}
\newcommand{\sC}{\mathbb{C}}
\newcommand{\sD}{\mathbb{D}}
\newcommand{\sE}{\mathbb{E}}
\newcommand{\sF}{\mathbb{F}}
\newcommand{\sG}{\mathbb{G}}
\newcommand{\sH}{\mathbb{H}}
\newcommand{\sI}{\mathbb{I}}
\newcommand{\sJ}{\mathbb{J}}
\newcommand{\sK}{\mathbb{K}}
\newcommand{\sL}{\mathbb{L}}
\newcommand{\sM}{\mathbb{M}}
\newcommand{\sN}{\mathbb{N}}
\newcommand{\sO}{\mathbb{O}}
\newcommand{\sP}{\mathbb{P}}
\newcommand{\sQ}{\mathbb{Q}}
\newcommand{\sR}{\mathbb{R}}
\newcommand{\sS}{\mathbb{S}}
\newcommand{\sT}{\mathbb{T}}
\newcommand{\sU}{\mathbb{U}}
\newcommand{\sV}{\mathbb{V}}
\newcommand{\sW}{\mathbb{W}}
\newcommand{\sX}{\mathbb{X}}
\newcommand{\sY}{\mathbb{Y}}
\newcommand{\sZ}{\mathbb{Z}}

% Frequently used operators.
\DeclareMathOperator{\Cov}{Cov}
\DeclareMathOperator{\Var}{Var}
\DeclareMathOperator{\diag}{diag}
\DeclareMathOperator{\nei}{Nei}
\DeclareMathOperator{\nnz}{nnz}
\DeclareMathOperator{\pr}{pr}
\DeclareMathOperator{\prox}{prox}
\DeclareMathOperator{\push}{push}
\DeclareMathOperator{\res}{res}
\DeclareMathOperator{\rel}{Rel}
\DeclareMathOperator{\relu}{ReLU}
\DeclareMathOperator{\sign}{sign}
\DeclareMathOperator{\supp}{supp}
\DeclareMathOperator{\tr}{tr}
\DeclareMathOperator{\Tr}{Tr}
\DeclareMathOperator{\vol}{vol}
\DeclareMathOperator{\work}{work}
\DeclareMathOperator{\Work}{Work}
\DeclareMathOperator*{\argmax}{arg\,max}
\DeclareMathOperator*{\argmin}{arg\,min}

% Project-wide names and claim-status labels used by standalone research notes.
% Mathematical notation belongs in math_commands.tex.

% Algorithms and solver families.
\newcommand{\AESP}{\textsc{AESP}}
\newcommand{\APPR}{\textsc{APPR}}
\newcommand{\ASPR}{\textsc{ASPR}}
\newcommand{\FISTA}{\textsc{FISTA}}
\newcommand{\ISTA}{\textsc{ISTA}}
\newcommand{\LOCAPPR}{\textsc{LocAPPR}}
\newcommand{\LOCGD}{\textsc{LocGD}}
\newcommand{\LOCSOR}{\textsc{LocSOR}}
\newcommand{\LOCCH}{\textsc{LocCH}}
\newcommand{\LOCHB}{\textsc{LocHB}}
\newcommand{\LocProxGS}{\textsc{LocProxGS}}
\newcommand{\Hybrid}{\textsc{AESP--LocSOR}}

% Claim classifications required by the repository research-note protocol.
\newcommand{\statussource}{\textbf{Source result}}
\newcommand{\statusproved}{\textbf{Proved here}}
\newcommand{\statusconditional}{\textbf{Conditional}}
\newcommand{\statusconjecture}{\textbf{Open / conjectural}}
\newcommand{\statusopen}{\textbf{Open}}
\newcommand{\statusempirical}{\textbf{Empirical}}
\newcommand{\statusobservation}{\textbf{Empirical observation}}
\newcommand{\statusrefuted}{\textbf{Refuted route}}

% Compatibility names used by the older complete-note presentation layer.
\newcommand{\Status}[2]{\textbf{#1:} #2}
\newcommand{\SourceStatus}{\textnormal{\textsc{Source result}}}
\newcommand{\ProvedStatus}{\textnormal{\textsc{Proved here}}}
\newcommand{\ConditionalStatus}{\textnormal{\textsc{Conditional}}}
\newcommand{\ComputedStatus}{\textnormal{\textsc{Computed}}}
\newcommand{\RefutedStatus}{\textnormal{\textsc{Refuted route}}}
\newcommand{\OpenStatus}{\textnormal{\textsc{Open}}}



\title{Sharp Post-Lock Warmup Certificates on Paths and Spiders}
\author{Baojian Zhou\\
\small Working research note for \texttt{hybrid-local-solver}}
\date{August 2026}

\begin{document}
\maketitle

\begin{abstract}
This note asks whether a componentwise entrance analysis can replace the
conservative spectral warmup used to prevent support retraction after a face
change.  The first target is an endpoint path and the second is a spider.
We first reconstruct the exact entrance condition in the project's
source-aligned variables.  An exact rational screen then separates two
questions that had been conflated: one stable-face proximal solve can be safe
on the realized path traces while still being insufficient uniformly over the
whole nonnegative residual cone.  In fact, supplied full paths give exact
cone witnesses against one and two proximal solves.  These witnesses do not
refute a constant warmup for the realized endpoint-seeded RPPR chronology;
that sharper question and the spider junction remain open.
\end{abstract}

% Canonical source-aligned PageRank/RPPR reference formulation.
%
% This fragment is intentionally heading-free at the section level so that the
% active paper and every standalone note can input it. It is a manuscript
% reference convention, not an implementation-wide residual decision.
% Primary source: Fountoulakis and Martinez-Rubio, arXiv:2602.21138v2,
% PDF p. 3, Section 3 and equation (RPPR).

\subsection{Common source-aligned PageRank and RPPR model}
\label{subsec:shared-source-aligned-problem}

All documents in this manuscript workspace use the following reference
language. It follows the plain italic notation of the comparison paper:
\(x,s,A,D,Q,I\) denote column vectors or matrices as determined by their
displayed domains. This is the scoped exception to the project's usual
bold-vector and bold-matrix typography. A note may impose a stronger
assumption after this block, but it must not redefine these objects.

Let \(G=(V,E)\) be a finite, simple, undirected, unweighted graph without
isolated vertices, let \(V=[n]:=\{1,\ldots,n\}\), and let
\(A\in\{0,1\}^{n\times n}\) be its symmetric adjacency matrix. For
\(i\in V\), write
\[
    \mathcal N(i):=\{j\in V:\{i,j\}\in E\},
    \qquad
    d_i:=|\mathcal N(i)|,
    \qquad
    D:=\diag(d_1,\ldots,d_n).
\]
For \(S\subseteq V\), define
\[
    \mathcal N(S):=\bigcup_{i\in S}\mathcal N(i),
    \qquad
    \partial S:=\mathcal N(S)\setminus S,
    \qquad
    \operatorname{Ext}(S):=V\setminus(S\cup\partial S),
\]
\begin{equation}
    \vol(S):=\sum_{i\in S}d_i.
    \label{eq:shared-volume}
\end{equation}
The support of a vector is
\(\supp(x):=\{i\in[n]:x_i\neq0\}\). Matrix inequalities use the Loewner
order, and all unqualified vector inequalities are coordinatewise.

The seed is a column distribution
\begin{equation}
    s\in\R_+^n,
    \qquad
    \inner{\one}{s}=1.
    \label{eq:shared-seed}
\end{equation}
The single-seed specialization is always written \(s=e_v\), where \(v\in V\)
is the seed vertex. Thus \(s\) is never used as a vertex label. For
\(\alpha\in(0,1]\), define
\begin{equation}
\begin{aligned}
    \mathcal L&:=I-D^{-1/2}AD^{-1/2},\\
    Q&:=\alpha I+\frac{1-\alpha}{2}\mathcal L
      =\frac{1+\alpha}{2}I
       -\frac{1-\alpha}{2}D^{-1/2}AD^{-1/2},\\
    b&:=\alpha D^{-1/2}s.
\end{aligned}
    \label{eq:shared-pagerank-matrices}
\end{equation}
Since \(0\preceq\mathcal L\preceq2I\),
\(\alpha I\preceq Q\preceq I\). The shared smooth PageRank quadratic is
\begin{equation}
    f(x):=\frac12\inner{x}{Qx}-\inner{b}{x},
    \qquad
    \nabla f(x)=Qx-b.
    \label{eq:shared-pagerank-objective}
\end{equation}
Its unique unregularized minimizer and associated lazy personalized PageRank
vector are
\begin{equation}
    x^0:=Q^{-1}b,
    \qquad
    \pi:=D^{1/2}x^0.
    \label{eq:shared-ppr-solution}
\end{equation}

For a regularization parameter \(\rho>0\), reserve \(g_\rho\) for the
nonsmooth weighted \(\ell_1\) term and define
\begin{equation}
    g_\rho(x):=\alpha\rho\norm{D^{1/2}x}_1,
    \qquad
    F_\rho(x):=f(x)+g_\rho(x),
    \qquad
    x^\star(\rho):=\argmin_{x\in\R^n}F_\rho(x).
    \label{eq:shared-rppr-objective}
\end{equation}
When \(\rho\) is fixed, \(x^\star\) abbreviates \(x^\star(\rho)\), never the
unregularized point \(x^0\). Write
\[
    S^\star(\rho):=\supp(x^\star(\rho)),
    \qquad
    I^\star(\rho):=[n]\setminus S^\star(\rho).
\]
The source records \(x^\star(\rho)\geq0\),
\(\vol(S^\star(\rho))\leq1/\rho\), and the coordinatewise conditions
\begin{equation}
    \nabla_i f(x^\star(\rho))
    \in
    \begin{cases}
        \{-\alpha\rho\sqrt{d_i}\}, & x_i^\star(\rho)>0,\\
        [-\alpha\rho\sqrt{d_i},0], & x_i^\star(\rho)=0.
    \end{cases}
    \label{eq:shared-rppr-kkt}
\end{equation}

The common computational unit is one repeated adjacency-list scan: scanning
coordinate \(i\) costs \(d_i\), and scanning a set \(S\) costs \(\vol(S)\).
Iteration count and wall-clock time do not replace this degree-weighted work
measure.

Finally, accuracy symbols remain distinct. The APPR threshold is
\(\epsappr\); \(\epsobj\) denotes an objective-gap target;
\(\epspg\) denotes the source experiment's proximal fixed-point tolerance; and
\(\epsppr\) may denote a document-scoped degree-normalized PPR error target.
No equality or conversion among these quantities is assumed. In particular,
this shared model does not resolve the repository-wide residual decision.


\section{Question and claim boundary}

The target is a proved constant or logarithmic post-lock warmup bound under a
fully stated changing-face recurrence and charged-work model.  A failed proof,
growing exact family, or spider obstruction will be retained as a refuted
route rather than promoted to an accelerated local-solver theorem.

The graph is a finite undirected endpoint path with ambient degrees, and all
variables below use the source-aligned PageRank parameter
\(q=\sqrt{\alpha/(1-\alpha)}\).  A supplied fixed face is not a local
discovery algorithm.  Every proximal solve, face sweep, inverse application,
certificate query, and state write must be charged by any algorithmic
corollary.  This note presently proves no \(\varepsilon_{\rm ppr}\) stopping
or end-to-end work theorem.

\subsection{Exact reconstruction of the Fable entrance gate}
\label{sec:gate-reconstruction}

Let \(\widetilde Q=D^{1/2}QD^{1/2}\),
\(\widetilde c=\alpha(s-\rho d)\), and
\(\kappa=1-2\alpha\), as in the exploratory Fable recurrence.  On a fixed
positive face \(S\), define the lower state residual
\begin{equation}
 Y_t=\widetilde c-\widetilde Qx_t
 \quad\hbox{and}\quad
 \mathcal M_S=\kappa D_S
 (\widetilde Q_{SS}+\kappa D_S)^{-1}.
 \label{eq:path-lock-hat-residual-map}
\end{equation}
If a stage has no retraction, exact obstacle optimality gives
\begin{equation}
 Y_{t+1}=\mathcal M_Su_t,
 \qquad
 u_t=(1+\beta)Y_t-\beta Y_{t-1}.
 \label{eq:path-lock-trigger-recurrence}
\end{equation}
The stage does not trigger if and only if \(u_t\geq0\) componentwise on its
support.  Conjugating by \(D^{1/2}\),
\begin{equation}
 D_S^{-1/2}\mathcal M_SD_S^{1/2}
 =\kappa(Q_{SS}+\kappa I)^{-1}.
 \label{eq:path-lock-normalized-map}
\end{equation}
Thus this is not the transported-center residual wave of the sibling
\texttt{path\_terminal\_modal\_block} note.  The two recurrences share the
weighted path cosine basis, but neither one's entry profile can be imported
into the other without a new calculation.

For completeness, split the normalized residual as
\(D^{-1/2}Y_t=L_t\phi+h_t\), where \(\phi>0\) is the Perron direction of the
normalized map.  Fable's sharp sufficient entrance test is
\begin{equation}
 \left|(1+\beta)h_t-\beta h_{t-1}\right|
 \leq \beta L_{t-1}\phi
 \quad\hbox{componentwise},
 \qquad
 L_t\geq(1-q_S)L_{t-1}>0.
 \label{eq:path-lock-sharp-entrance}
\end{equation}
The absolute values make
\eqref{eq:path-lock-sharp-entrance} stronger than the exact trigger
\(u_t\geq0\).  The conservative certified warmup controls it by one global
norm; the present target is the true componentwise condition.

\subsection{One prox solve certifies the first momentum stage}
\label{sec:first-momentum-safe}

There is a useful graph-uniform fact hidden by the conservative modal norm.
It certifies entrance to momentum, but not the whole momentum tail.

\begin{proposition}[First-momentum entrance; proved here]
\label{prop:first-momentum-safe}
Let \(0<\alpha\leq1/5\), let \(S\) be any supplied face, and let
\(0\leq\beta\leq1\).  Suppose a stable-face pure proximal solve starts from
a state residual \(Y\geq0\) and produces \(Y^+=\mathcal M_SY\).  Then the
immediately following momentum trigger is nonnegative:
\begin{equation}
 (1+\beta)Y^+-\beta Y\geq0
 \quad\hbox{componentwise}.
 \label{eq:first-momentum-safe}
\end{equation}
Consequently that first momentum stage has no retraction.  The claim applies
to every endpoint-path prefix and every spider face, without a spectral-gap
or Perron-spread constant.
\end{proposition}

\begin{proof}
Write \(S_S=D_S^{-1/2}A_{SS}D_S^{-1/2}\geq0\).  Direct substitution gives
\begin{equation}
 \overline{\mathcal M}_S
 :=D_S^{-1/2}\mathcal M_SD_S^{1/2}
 =\frac{2(1-2\alpha)}{1-\alpha}(3I-S_S)^{-1}.
 \label{eq:first-momentum-neumann-map}
\end{equation}
Since \(\lVert S_S\rVert_2\leq1\), the Neumann series
\begin{equation}
 (3I-S_S)^{-1}=\frac13\sum_{r\geq0}(S_S/3)^r
 \label{eq:first-momentum-neumann-series}
\end{equation}
is entrywise nonnegative and has diagonal at least \(1/3\).  The scalar
coefficient in \eqref{eq:first-momentum-neumann-map} is at least \(3/2\)
exactly when \(\alpha\leq1/5\).  Hence
\(\overline{\mathcal M}_S\geq\frac12I\) entrywise.  With
\(\sigma=\beta/(1+\beta)\leq1/2\), this gives
\(\overline{\mathcal M}_S-\sigma I\geq0\).  Positive diagonal conjugation
preserves entrywise signs, so
\[
 (1+\beta)Y^+-\beta Y
 =(1+\beta)(\mathcal M_S-\sigma I)Y\geq0.
\]
The exact trigger equivalence in
\eqref{eq:path-lock-trigger-recurrence} completes the proof.
\end{proof}

This proposition reduces the actual entrance warmup from Fable's large
norm-certified \(J_S\) to one solve if the goal is only to launch one safe
momentum stage.  It does not establish
\eqref{eq:path-lock-sharp-entrance}, and
\cref{prop:path-lock-cone-witnesses} shows why the conclusion cannot simply
be iterated forever.

\begin{proposition}[Alternating reset loses acceleration; proved here]
\label{prop:alternating-reset-slow}
Consider the full connected face and repeat one pure proximal stage followed
by one momentum stage, resetting momentum after every such pair.  On the
Perron residual mode, one two-stage cycle has exact multiplier
\begin{equation}
 \eta_q=(1-q)^2(1+2q)=1-3q^2+2q^3.
 \label{eq:alternating-reset-multiplier}
\end{equation}
Thus for \(0<q\leq1/2\), a constant-factor reduction of a nonzero Perron
component needs \(\Theta(q^{-2})=\Theta(\alpha^{-1})\) cycles.  Alternating
the proved safe entrance with a reset is not an accelerated algorithm.
\end{proposition}

\begin{proof}
On the full face the Perron eigenvalue of \(\mathcal M\) is
\(m_0=1-q^2\), and the global momentum is
\(\beta=(1-q)/(1+q)\).  One prox--momentum pair maps a residual by
\(\mathcal M((1+\beta)\mathcal M-\beta I)\).  Its Perron multiplier is
\[
 m_0((1+\beta)m_0-\beta)=(1-q)^2(1+2q).
\]
Finally, \(1-\eta_q=q^2(3-2q)\) lies between \(2q^2\) and \(3q^2\) on the
stated interval.
\end{proof}

\subsection{A cone-uniform shortcut is already false}
\label{sec:cone-screen}

Suppose a supplied full face receives an arbitrary nonnegative residual
column \(y\).  After \(J\geq1\) stable-face pure proximal solves, the two
entrance residuals are \(\mathcal M^{J-1}y\) and
\(\mathcal M^Jy\).  The first trigger matrix is therefore
\begin{equation}
 U_0^{(J)}=
 \mathcal M^{J-1}\bigl((1+\beta)\mathcal M-\beta I\bigr),
 \label{eq:path-lock-trigger-kernel-initial}
\end{equation}
and subsequent trigger matrices obey
\begin{equation}
 U_{k+1}^{(J)}=
 \mathcal M\bigl((1+\beta)U_k^{(J)}-\beta U_{k-1}^{(J)}\bigr),
 \qquad U_{-1}^{(J)}=\mathcal M^{J-1}.
 \label{eq:path-lock-trigger-kernel-recurrence}
\end{equation}

\begin{proposition}[Exact cone witnesses; proved here]
\label{prop:path-lock-cone-witnesses}
A theorem asserting permanent no-triggering after one stable-face proximal
solve for \emph{every} nonnegative residual is false, even on the full
endpoint path.  On \(P_4\), with \(q=1/8\), \(J=1\), the endpoint diagonal
of \(U_2^{(1)}\) is
\begin{equation}
 (U_2^{(1)})_{00}=-\frac{541}{48000}<0.
 \label{eq:path-lock-one-warmup-witness}
\end{equation}
The analogous two-solve assertion is also false: on \(P_8\), with
\(q=1/16\), \(J=2\),
\begin{equation}
 (U_3^{(2)})_{00}
 =-\frac{151523312832478779015383634375}
 {28857779215644070278768053190656}<0.
 \label{eq:path-lock-two-warmup-witness}
\end{equation}
In each case the nonnegative entrance column \(y=e_0\) realizes the negative
trigger coordinate.

Three solves are not cone-uniform either.  On \(P_{46}\), with
\(q=1/92\), exact rational elimination gives
\begin{equation}
 (U_4^{(3)})_{00}<0.
 \label{eq:path-lock-three-warmup-witness}
\end{equation}
The verifier prints the full reduced fraction; its decimal value is about
\(-1.0542862762\times10^{-5}\).  The smaller sizes through \(P_{44}\) do not
replace this witness: finite endpoint effects delay the first negative
diagonal until size 46 in the screened family.
\end{proposition}

\begin{proof}
For a full endpoint path all entries of \(\widetilde Q\), \(D\),
\(\mathcal M\), and \(\beta=(1-q)/(1+q)\) are rational.  Exact Gaussian
elimination followed by
\eqref{eq:path-lock-trigger-kernel-initial} and
\eqref{eq:path-lock-trigger-kernel-recurrence}
gives the two displayed fractions and the exact negative sign in
\eqref{eq:path-lock-three-warmup-witness}.  Multiplication by \(e_0\)
selects the first column, so its endpoint is negative and the exact trigger
condition fails.
\end{proof}

The witnesses are deliberately scoped.  They refute a proof that forgets the
actual face-entry profile and uses only \(Y\geq0\); they do not show that
\(e_0\) is reached by the endpoint-seeded changing-face RPPR execution.  They
also do not rule out a different absolute constant.  The exact finite screen
over \(q=1/(2n)\) found the first horizon-passing values \(J=2\) on
\(P_4,P_6\) and \(J=3\) on \(P_8,P_{10},P_{12},P_{14},P_{16},P_{20}\),
through 32 or 64 momentum stages as applicable.  This passing table is
\emph{Measured} finite exact arithmetic, not an all-time theorem.

\section{Open path and spider targets}
\label{sec:path-spider-targets}

The first live path target is now sharper: characterize the residual produced
by one stable-face proximal solve after an actual endpoint-prefix admission,
then prove or refute that
\eqref{eq:path-lock-sharp-entrance} or the weaker exact trigger remains valid
for the whole segment after a graph-size-independent number of solves.  A
constant must be proved on the realized profile, not on the full nonnegative
cone unless the latter is genuinely true.

Only after that path statement survives should it be extended to a spider.
In normalized variables each arm carries the same one-dimensional resolvent
recurrence, while the center contributes one exact coupling equation whose
degree is the number of arms.  A spider theorem must state whether the arm
count and unequal arm lengths enter the warmup, and must charge a center scan,
all touched arm rows, and every repeated solve.  No such theorem is claimed
yet.

\section{Reproducibility}

The note-local script \texttt{verify\_warmup.py} constructs
\(\mathcal M\) and every trigger matrix with exact rational arithmetic.  Run
\begin{verbatim}
python3 verify_warmup.py --sizes 4 6 8 10 12 --horizon 32
python3 verify_warmup.py --sizes 14 16 20 --horizon 64
python3 verify_warmup.py --sizes 46 --horizon 4 --max-warmup 3
\end{verbatim}
The script's first negative entry is a proof witness.  A horizon pass remains
a finite screen and is printed with that name.

\end{document}
